\section{Changes of ownership in encoded frames}
\label{sec:peps_ownership_changes}
The statements below formalize Lemmas~6.5 and~6.6 of
\cite{OpenAI2026PolynomialPEPS}, September 24, 2026, \texttt{05-frames.tex},
lines 347--570, for the canonical maps constructed in their proofs, with the
information bounds kept as hypotheses, together with their reading as bounded
changes: allowed monomials of \cite[Theorem~5.2]{OpenAI2026PolynomialPEPS}
(Section~\ref{sec:peps_pair_effects}) on the registers of the frames, one per
site and per tag (Definition~\ref{def:peps_frame_registers}). The monomials are
first constructed on registers in which the raw registers of a region held by
one party are grouped into one register, and then placed on the registers of
the frames by reorderings of tensor factors and grouping unitaries.

We keep the setting of Section~\ref{sec:peps_encoded_frames}. For a set
$S\subseteq\Lambda$ write $\mathcal H_S=\bigotimes_{x\in S}\C^q$. For a vector
$\Omega\in\mathcal H_\Lambda$ and disjoint sets $T,E\subseteq\Lambda$, the
mutual information $I_\Omega(T:E)$ is
$S(\rho_T)+S(\rho_E)-S(\rho_{TE})$ with natural logarithms, where the $\rho$
are the reduced density matrices of $\ket\Omega\bra\Omega$.

\subsection{Splitting a sheet}

\begin{lemma}[Operators on parts of a sheet]
    \label{lem:peps_frame_split_forms}
    \lean{TNLean.PEPS.EncodedFrame.sheetSplit,
        TNLean.PEPS.EncodedFrame.threeSplit,
        TNLean.PEPS.EncodedFrame.rectKronecker_submatrix_sheetSplit,
        TNLean.PEPS.EncodedFrame.rectKronecker_submatrix_threeSplit,
        TNLean.PEPS.EncodedFrame.exists_kronecker_one_of_mem_supportedOperators,
        TNLean.PEPS.EncodedFrame.exists_one_kronecker_of_mem_supportedOperators,
        TNLean.PEPS.EncodedFrame.exists_threeSplit_eq_kronecker_one,
        TNLean.PEPS.EncodedFrame.exists_threeSplit_eq_one_kronecker_kronecker_one}
    \leanok
    Identify $\mathcal H_\Lambda$ with $\mathcal H_S\otimes\mathcal
        H_{\Lambda\setminus S}$, and, for disjoint $T$ and $E$ with
    $U=\Lambda\setminus(T\cup E)$, with $\mathcal H_T\otimes\mathcal H_E\otimes
        \mathcal H_U$. An operator acting on the sites of $S$ is $A'\otimes\id$, an
    operator acting on a set disjoint from $S$ is $\id\otimes A'$, an operator
    acting on $T\cup E$ is $A'\otimes\id_U$, and an operator acting on a subset
    of $E$ is $\id_T\otimes A'\otimes\id_U$.
\end{lemma}

\begin{proof}\leanok
    A product operator $\bigotimes_xm_x$ is, in these coordinates, the tensor
    product of its restrictions to the parts, and $m_x=\id$ off the support.
    Both sides are linear in the operator.
\end{proof}

\begin{lemma}[Splitting consequence on a sheet]
    \label{lem:peps_frame_splitting}
    \lean{TNLean.PEPS.EncodedFrame.SplittingData,
        TNLean.PEPS.EncodedFrame.SplittingData.error,
        TNLean.PEPS.EncodedFrame.splitVec,
        TNLean.PEPS.EncodedFrame.partialTraceRight_splitVec,
        TNLean.PEPS.EncodedFrame.quantumRelativeEntropy_splitVec_eq_mutualInformation,
        TNLean.PEPS.EncodedFrame.exists_sheetSplitting,
        TNLean.PEPS.EncodedFrame.exists_sheetSplitting_zpow}
    \leanok
    \uses{lem:peps_frame_split_forms}
    Let $T,E\subseteq\Lambda$ be disjoint, $U=\Lambda\setminus(T\cup E)$,
    $\lVert\Omega\rVert=1$ and $b=I_\Omega(T:E)$. There are finite-dimensional
    spaces $\mathcal B_T,\mathcal B_E$, an isometry
    $V\colon\mathcal H_U\to\mathcal B_T\otimes\mathcal B_E$ and unit vectors
    $s\in\mathcal H_T\otimes\mathcal B_T$, $s'\in\mathcal H_E\otimes\mathcal
        B_E$ with
    \begin{align}
        \lVert(\id_{TE}\otimes V)\Omega-s\otimes s'\rVert
        \le\sqrt{2(1-e^{-b/2})}\le\sqrt b .
    \end{align}
    In particular, if $b\le L^{-60}$ the error is at most $L^{-30}$. This is
    \cite[Lemma~6.4]{OpenAI2026PolynomialPEPS} (\texttt{05-frames.tex},
    lines 352--368) on the sheet.
\end{lemma}

\begin{proof}\leanok
    The relative entropy of $\rho_{TE}$ with respect to $\rho_T\otimes\rho_E$
    is $b$. The splitting estimate for a unit vector on three systems, with
    $\mathcal B_T=\mathcal H_T$ and $\mathcal B_E=\mathcal H_E\oplus\mathcal
        H_U$, applies in the coordinates $\mathcal H_T\otimes\mathcal
        H_E\otimes\mathcal H_U$.
\end{proof}

\subsection{Homogeneous birth and death}

\begin{definition}[Canonical birth map]
    \label{def:peps_frame_birth_map}
    \lean{TNLean.PEPS.EncodedFrame.pairSource,
        TNLean.PEPS.EncodedFrame.birthKernel,
        TNLean.PEPS.EncodedFrame.birthEffect,
        TNLean.PEPS.EncodedFrame.birthOp,
        TNLean.PEPS.EncodedFrame.sheetBirthOp,
        TNLean.PEPS.EncodedFrame.birthOp_eq,
        TNLean.PEPS.EncodedFrame.pairSource_mul_conjTranspose_apply}
    \leanok
    \uses{lem:peps_frame_splitting}
    With the data of the splitting lemma put $\mathcal V=\id_{TE}\otimes V$ and
    let $S\colon\mathcal H_E\otimes\mathcal B_E\to\mathcal
        H_{TE}\otimes\mathcal B_T\otimes\mathcal B_E$ insert $s$ on
    $T\mathcal B_T$. The canonical map of the birth is
    \begin{align}
        B=\mathcal V^\dagger SS^\dagger\mathcal V
        =\mathcal V^\dagger\bigl(\ket s\bra s_{T\mathcal B_T}
        \otimes\id_{E\mathcal B_E}\bigr)\mathcal V
    \end{align}
    (\texttt{05-frames.tex}, equation (birth-map), lines 414--426). It is the
    effect $S^\dagger\mathcal V$, which applies $V$ on $U$ and contracts
    $T\mathcal B_T$ with $\bra s$, followed by its adjoint, which supplies
    $s$ and applies $V^\dagger$.
\end{definition}

\begin{lemma}[Properties of the birth map]
    \label{lem:peps_frame_birth_map}
    \lean{TNLean.PEPS.EncodedFrame.pairSource_conjTranspose_mul_self,
        TNLean.PEPS.EncodedFrame.pairSource_mulVec,
        TNLean.PEPS.EncodedFrame.birthOp_isHermitian,
        TNLean.PEPS.EncodedFrame.norm_birthOp_le_one,
        TNLean.PEPS.EncodedFrame.norm_act_birthOp_sub_le,
        TNLean.PEPS.EncodedFrame.birthEffect_eq_submatrix,
        TNLean.PEPS.EncodedFrame.birthOp_eq_submatrix,
        TNLean.PEPS.EncodedFrame.commute_birthOp,
        TNLean.PEPS.EncodedFrame.sheetBirthOp_isHermitian,
        TNLean.PEPS.EncodedFrame.norm_sheetBirthOp_le_one,
        TNLean.PEPS.EncodedFrame.norm_act_sheetBirthOp_sub_le,
        TNLean.PEPS.EncodedFrame.commute_sheetBirthOp}
    \leanok
    \uses{def:peps_frame_birth_map}
    The map $B$ is a Hermitian contraction, it acts as the identity on
    $\mathcal H_E$, so it commutes with every operator acting on a subset of
    $E$, and for every vector $\psi$
    \begin{align}
        \lVert(B-\id)\psi\rVert\le\lVert\mathcal V\psi-s\otimes s'\rVert .
    \end{align}
\end{lemma}

\begin{proof}\leanok
    Since $s$ is a unit vector, $S^\dagger S=\id$, so $\Pi=SS^\dagger$ is an
    orthogonal projection with $\Pi(s\otimes s')=s\otimes s'$. As
    $\mathcal V^\dagger\mathcal V=\id$,
    $(B-\id)\psi=\mathcal V^\dagger(\Pi-\id)(\mathcal V\psi-s\otimes s')$, and
    $\id-\Pi$ is a projection. In the coordinates
    $\mathcal H_E\otimes(\mathcal H_T\otimes\mathcal H_U)$ the effect
    $S^\dagger\mathcal V$ is $\id_E\otimes K$ with $K=\bra s(\id_T\otimes V)$,
    so $B=\id_E\otimes K^\dagger K$.
\end{proof}

\begin{lemma}[Operators passing through an encoding]
    \label{lem:peps_frame_pass_encoding}
    \lean{TNLean.PEPS.EncodedFrame.one_kronecker_mul_stack,
        TNLean.PEPS.EncodedFrame.one_kronecker_mul_frameEncoder,
        TNLean.PEPS.EncodedFrame.norm_act_one_kronecker_refVec_sub_le}
    \leanok
    \uses{def:peps_frame}
    If a raw operator $B$ commutes with the raw part of the encoding of a
    frame on every tag configuration, then $(\id_{\rm tags}\otimes
        B)K_{\mathcal F}=K_{\mathcal F}B$ and
    $\lVert(\id\otimes B)\Omega_{\mathcal F}-\Omega_{\mathcal F}\rVert
        \le\lVert B\Omega-\Omega\rVert$.
\end{lemma}

\begin{proof}\leanok
    Compare the two sides on each tag; the encoding is a contraction.
\end{proof}

\begin{definition}[Birth partition]
    \label{def:peps_frame_birth_partition}
    \lean{TNLean.PEPS.EncodedFrame.Frame.changeOwner,
        TNLean.PEPS.EncodedFrame.Frame.birthEnv,
        TNLean.PEPS.EncodedFrame.Frame.mem_birthEnv,
        TNLean.PEPS.EncodedFrame.Frame.outerHoles_subset_birthEnv}
    \leanok
    \uses{def:peps_frame}
    For an encoded frame $\mathcal F$, a party $P_\circ$ and a set $T$ of raw
    sites, let $\mathcal F'$ be the frame in which the raw registers of $T$
    are owned by $Q_\circ$, with every hole, encoding and tag owner kept, and
    \begin{align}
        E=\{x:\text{the raw owner of }x\neq P_\circ\}\cup H^+(\mathcal F),
        \qquad U=\Lambda\setminus(T\cup E)
    \end{align}
    (\texttt{05-frames.tex}, equation (birth-partition), lines 400--405).
\end{definition}

\begin{lemma}[Owners in the birth partition]
    \label{lem:peps_frame_birth_partition_owners}
    \lean{TNLean.PEPS.EncodedFrame.Frame.owner_eq_of_notMem_birthEnv,
        TNLean.PEPS.EncodedFrame.Frame.changeOwner_owner_of_notMem_birthEnv,
        TNLean.PEPS.EncodedFrame.Frame.disjoint_birthEnv}
    \leanok
    \uses{def:peps_frame_birth_partition}
    Every register outside $E$ is owned by $P_\circ$ in $\mathcal F$, and by
    $P_\circ$ or $Q_\circ$ in $\mathcal F'$. If $T$ is owned by $P_\circ$ and
    avoids every outer hole footprint, then $T$ and $E$ are disjoint
    (\texttt{05-frames.tex}, lines 397--405).
\end{lemma}

\begin{proof}\leanok
    A register outside $E$ has old owner $P_\circ$ by the definition of $E$;
    in $\mathcal F'$ its owner is $Q_\circ$ if it lies in $T$ and its old
    owner otherwise. A register of $T$ has old owner $P_\circ$ and lies in no
    outer hole footprint, so it is in neither part of $E$.
\end{proof}

\begin{theorem}[Homogeneous birth and death: canonical map]
    \label{thm:peps_frame_birth}
    \lean{TNLean.PEPS.EncodedFrame.Frame.birth,
        TNLean.PEPS.EncodedFrame.Frame.death}
    \leanok
    \uses{lem:peps_frame_birth_map,lem:peps_frame_pass_encoding,
        def:peps_frame_birth_partition,lem:peps_frame_birth_partition_owners}
    In an encoded frame $\mathcal F$, let $T$ be owned by $P_\circ$ and
    disjoint from every outer hole footprint, let $E$ be as in the birth
    partition, and let $\lVert\Omega\rVert=1$. If $I_\Omega(T:E)\le L^{-60}$,
    the canonical map $B$ of splitting data with error at most $L^{-30}$ is a
    Hermitian contraction, acts as the identity on $E$ and on all tags,
    satisfies $(\id\otimes B)K_{\mathcal F}=K_{\mathcal F'}B$, and
    \begin{align}
        \lVert(\id\otimes B)\Omega_{\mathcal F}-\Omega_{\mathcal F'}\rVert
        \le L^{-30}.
    \end{align}
    Since $\mathcal F$ and $\mathcal F'$ differ only in raw owners,
    $K_{\mathcal F'}=K_{\mathcal F}$ and $\Omega_{\mathcal
            F'}=\Omega_{\mathcal F}$, so the intertwining identity says that $B$
    commutes with the encoding; the change of owner is the content of the
    bounded-change assertion below. The death is the same map with the
    reverse input and output ownership; its conditions are those of the frame
    after the death. This is the
    content of \cite[Lemma~6.5]{OpenAI2026PolynomialPEPS} for the canonical
    map, with its hypotheses.
\end{theorem}

\begin{proof}\leanok
    \uses{lem:peps_frame_splitting,lem:peps_frame_birth_map,
        lem:peps_frame_pass_encoding,def:peps_frame_birth_partition,
        lem:peps_frame_birth_partition_owners}
    Every hole encoder acts on $H^+(\mathcal F)\subseteq E$, so it commutes
    with $B$. Hence $\lVert(\id\otimes B)\Omega_{\mathcal F}-\Omega_{\mathcal
            F}\rVert\le\lVert(B-\id)\Omega\rVert\le\lVert\mathcal V\Omega-s\otimes
        s'\rVert\le L^{-30}$, and $\Omega_{\mathcal F'}=\Omega_{\mathcal F}$ in
    canonical coordinates, since only raw owners change.
\end{proof}

\begin{definition}[Parties used by a monomial]
    \label{def:peps_monomial_parties}
    \lean{TNLean.PEPS.PairEffect.Word.UsesOnly,
        TNLean.PEPS.PairEffect.Word.sourceCount,
        TNLean.PEPS.PairEffect.PartyChain.UsesOnly,
        TNLean.PEPS.PairEffect.PartyChain.sourceCount}
    \leanok
    \uses{def:peps_party_layout,def:peps_party_monomial}
    An allowed monomial on parties uses only the parties of a set $S$ when
    each of its local maps acts at a party of $S$ and each of its pair sources
    and pair effects joins two parties of $S$; reorderings of tensor factors
    use no party. Its numbers of pair sources and of pair effects are counted
    along the composition.
\end{definition}

\begin{lemma}[Matrices on registers of one party]
    \label{lem:peps_front_registers}
    \lean{TNLean.PEPS.PairEffect.matL,
        TNLean.PEPS.PairEffect.matL_apply,
        TNLean.PEPS.PairEffect.ofLp_matL,
        TNLean.PEPS.PairEffect.norm_matL,
        TNLean.PEPS.PairEffect.matL_add,
        TNLean.PEPS.PairEffect.matL_mul,
        TNLean.PEPS.PairEffect.oneIso,
        TNLean.PEPS.PairEffect.twoIso,
        TNLean.PEPS.PairEffect.pairHeadIso_eq,
        TNLean.PEPS.PairEffect.appendIso_one_symm,
        TNLean.PEPS.PairEffect.pairHeadIso_appendIso_symm,
        TNLean.PEPS.PairEffect.matLocal₂₁,
        TNLean.PEPS.PairEffect.matLocal₁₂,
        TNLean.PEPS.PairEffect.matLocal₂₂,
        TNLean.PEPS.PairEffect.eval_localMap,
        TNLean.PEPS.PairEffect.eval_localMap₂₁,
        TNLean.PEPS.PairEffect.eval_localMap₁₂,
        TNLean.PEPS.PairEffect.pairHeadIso_eval_localMap₂₂}
    \leanok
    \uses{def:peps_party_layout}
    Identify one register $\C^\alpha$ with $\C^\alpha$, and two registers
    $\C^\alpha$, $\C^\beta$ with $\C^{\alpha\times\beta}$. A matrix $A$ from
    $\C^{\alpha\times\beta}$ (or $\C^\alpha$) to $\C^{\gamma\times\delta}$ (or
    $\C^\gamma$) defines a local map of a party $p$ on registers of $p$ of these
    dimensions, of norm at most $\lVert A\rVert$. Placed in front of an
    arbitrary list of further registers, its operator is $A\otimes\id$ in
    these identifications.
\end{lemma}

\begin{proof}\leanok
    Check the identity on pure tensors, using linearity.
\end{proof}

\begin{definition}[Birth and death monomials]
    \label{def:peps_frame_birth_monomials}
    \lean{TNLean.PEPS.EncodedFrame.birthWord,
        TNLean.PEPS.EncodedFrame.birthPrivateWord,
        TNLean.PEPS.EncodedFrame.deathChain,
        TNLean.PEPS.EncodedFrame.deathPrivateChain}
    \leanok
    \uses{def:peps_frame_birth_map,def:peps_party_monomial,lem:peps_front_registers}
    Let $K=\bra s(\id_T\otimes V)\colon\mathcal H_T\otimes\mathcal H_U\to\mathcal
        B_E$, with the data of the splitting lemma. Consider a register
    $\mathcal H_U$ of $P_\circ$ and further registers left untouched. For
    $Q_\circ\neq P_\circ$, the birth monomial starts from a register $\mathcal
        H_T$ of $P_\circ$: it applies the private contraction $K$ at $P_\circ$ from
    $\mathcal H_T\otimes\mathcal H_U$ to a register $\mathcal B_E$ of $P_\circ$,
    prepares the normalized pair vector $s$ on a new register $\mathcal H_T$ of
    $Q_\circ$ and a new register $\mathcal B_T$ of $P_\circ$, and applies the
    private contraction $V^\dagger$ at $P_\circ$ from $\mathcal
        B_T\otimes\mathcal B_E$ to $\mathcal H_U$ (\texttt{05-frames.tex}, lines
    413--419). The death monomial starts from a register $\mathcal H_T$ of
    $Q_\circ$: it applies the private isometry $V$ at $P_\circ$ from $\mathcal
        H_U$ to $\mathcal B_T\otimes\mathcal B_E$, contracts $\mathcal H_T$ and
    $\mathcal B_T$ with the normalized pair effect $\bra s$, and applies the
    private contraction $K^\dagger$ at $P_\circ$ from $\mathcal B_E$ to new
    registers $\mathcal H_T$, $\mathcal H_U$ of $P_\circ$ (lines 439--446). For
    $Q_\circ=P_\circ$ both are the private contraction $K^\dagger K$ at
    $P_\circ$ (lines 447--449).
\end{definition}

\begin{lemma}[Operators of the birth and death monomials]
    \label{lem:peps_frame_birth_monomial_operators}
    \lean{TNLean.PEPS.EncodedFrame.pairHeadIso_birthWord_eval,
        TNLean.PEPS.EncodedFrame.pairHeadIso_deathChain_eval,
        TNLean.PEPS.EncodedFrame.pairHeadIso_birthPrivateWord_eval,
        TNLean.PEPS.EncodedFrame.sourceThenAdjoint,
        TNLean.PEPS.EncodedFrame.sourceThenAdjoint_pairVec,
        TNLean.PEPS.EncodedFrame.pairHeadIso_frame_localMap_source,
        TNLean.PEPS.EncodedFrame.effectNext,
        TNLean.PEPS.EncodedFrame.effectNext_apply,
        TNLean.PEPS.EncodedFrame.effectNext_matL,
        TNLean.PEPS.EncodedFrame.norm_birthKernel_le_one,
        Matrix.IsIsometry.l2_opNorm_le_one,
        TNLean.PEPS.EncodedFrame.norm_toLp_eq_one_of_star_dotProduct,
        TNLean.PEPS.EncodedFrame.isAllowed_birthWord,
        TNLean.PEPS.EncodedFrame.isAllowed_birthPrivateWord,
        TNLean.PEPS.EncodedFrame.isAllowed_deathChain,
        TNLean.PEPS.EncodedFrame.usesOnly_birthWord,
        TNLean.PEPS.EncodedFrame.sourceCount_birthWord,
        TNLean.PEPS.EncodedFrame.usesOnly_birthPrivateWord,
        TNLean.PEPS.EncodedFrame.sourceCount_birthPrivateWord,
        TNLean.PEPS.EncodedFrame.usesOnly_deathChain,
        TNLean.PEPS.EncodedFrame.sourceCount_deathChain,
        TNLean.PEPS.EncodedFrame.effectCount_deathChain,
        TNLean.PEPS.EncodedFrame.sheetBirthOp_eq_submatrix}
    \leanok
    \uses{def:peps_frame_birth_monomials,def:peps_monomial_parties,
        lem:peps_frame_birth_map}
    The birth and death monomials are allowed and use only $P_\circ$ and
    $Q_\circ$; for $Q_\circ\neq P_\circ$ the birth has one pair source and no
    pair effect, and the death one pair effect and no pair source. With the
    registers $\mathcal H_T$ and $\mathcal H_U$ identified with $\mathcal
        H_T\otimes\mathcal H_U$, the operator of each is $K^\dagger K\otimes\id$ on
    the untouched registers, and on the sheet $B=\id_E\otimes K^\dagger K$.
\end{lemma}

\begin{proof}\leanok
    The contraction $K$ is the effect $S^\dagger\mathcal V$ with $E$ trivial, so
    $\lVert K\rVert\le1$, and $\lVert V\rVert\le1$. For the birth, on $t\otimes
        u\otimes w$ the first contraction gives $k=K(t\otimes u)$; preparing
    $s=\sum t'\otimes b$ and applying $V^\dagger$ to $b\otimes k$ gives
    $\sum t'\otimes V^\dagger(b\otimes k)=K^\dagger k$, by comparing
    coordinates and using linearity in the prepared vector. For the death,
    $V u=\sum b\otimes e$ and the effect gives $\sum\langle s,t\otimes
        b\rangle e=K(t\otimes u)$, after which $K^\dagger$ is applied.
\end{proof}

\begin{theorem}[Homogeneous birth and death on grouped registers]
    \label{thm:peps_frame_birth_bounded_grouped}
    \lean{TNLean.PEPS.EncodedFrame.Frame.birth_monomial,
        TNLean.PEPS.EncodedFrame.Frame.death_monomial,
        TNLean.PEPS.EncodedFrame.Frame.exists_birthWord,
        TNLean.PEPS.EncodedFrame.Frame.exists_deathChain}
    \leanok
    \uses{thm:peps_frame_birth,lem:peps_frame_birth_partition_owners,
        lem:peps_frame_birth_monomial_operators}
    In the setting of the canonical birth map, with its hypotheses, there are
    splitting data with the reference error of Theorem~\ref{thm:peps_frame_birth}
    such that $B=\id_E\otimes K^\dagger K$ with $K=\bra s(\id_T\otimes V)$, the
    registers of $U$ are owned by $P_\circ$ before and after the birth, and
    those of $T$ by $P_\circ$ before and by $Q_\circ$ after it. Group the raw
    registers of $T$ into one register $\mathcal H_T$ and those of $U$ into one
    register $\mathcal H_U$, followed by an arbitrary list of further registers.
    On these registers there is an allowed monomial using only $P_\circ$ and
    $Q_\circ$, with at most one normalized pair source, from $\mathcal H_T$ held
    by $P_\circ$ to $\mathcal H_T$ held by $Q_\circ$, whose operator is
    $K^\dagger K\otimes\id$; for $Q_\circ\neq P_\circ$ it is a private
    contraction at $P_\circ$, one pair source between $Q_\circ$ and $P_\circ$,
    and a private contraction at $P_\circ$. The death is the same with no pair
    source and at most one normalized pair effect. These grouped monomials are
    the middle steps of the bounded changes on the registers of the frame in
    Theorem~\ref{thm:peps_frame_birth_bounded}.
\end{theorem}

\begin{proof}\leanok
    \uses{thm:peps_frame_birth,lem:peps_frame_birth_monomial_operators,
        lem:peps_frame_birth_partition_owners}
    Take the splitting data of Theorem~\ref{thm:peps_frame_birth}. The
    registers of $U$ lie outside $E$, so they are owned by $P_\circ$ before
    and after, and those of $T$ are owned by $P_\circ$ before and by $Q_\circ$
    after. The birth and death monomials have operator $K^\dagger K$ on
    $\mathcal H_T\otimes\mathcal H_U$, and $B=\id_E\otimes K^\dagger K$.
\end{proof}

\begin{lemma}[Reordering registers]
    \label{lem:peps_register_reorderings}
    \lean{TNLean.PEPS.EncodedFrame.relabelHead,
        TNLean.PEPS.EncodedFrame.relabelSites,
        TNLean.PEPS.EncodedFrame.eval_relabelSites,
        TNLean.PEPS.EncodedFrame.partSites,
        TNLean.PEPS.EncodedFrame.partSites_eq,
        TNLean.PEPS.EncodedFrame.partWord,
        TNLean.PEPS.EncodedFrame.unpartWord,
        TNLean.PEPS.EncodedFrame.eval_partWord,
        TNLean.PEPS.EncodedFrame.eval_unpartWord,
        TNLean.PEPS.EncodedFrame.partWord_props,
        TNLean.PEPS.EncodedFrame.unpartWord_props,
        TNLean.PEPS.PairEffect.Word.usesOnly_moveHead,
        TNLean.PEPS.PairEffect.Word.sourceCount_moveHead,
        TNLean.PEPS.PairEffect.Word.usesOnly_exchangeBlocks,
        TNLean.PEPS.PairEffect.Word.sourceCount_exchangeBlocks,
        TNLean.PEPS.PairEffect.Word.eval_frameList_appendIso_symm,
        TNLean.PEPS.PairEffect.Word.assocWord,
        TNLean.PEPS.PairEffect.Word.unassocWord,
        TNLean.PEPS.PairEffect.Word.eval_assocWord_appendIso_symm,
        TNLean.PEPS.PairEffect.Word.eval_unassocWord_appendIso_symm,
        TNLean.PEPS.PairEffect.Word.swapPairs,
        TNLean.PEPS.PairEffect.Word.eval_swapPairs,
        TNLean.PEPS.EncodedFrame.partWordApp,
        TNLean.PEPS.EncodedFrame.unpartWordApp,
        TNLean.PEPS.EncodedFrame.eval_partWordApp,
        TNLean.PEPS.EncodedFrame.eval_unpartWordApp,
        TNLean.PEPS.EncodedFrame.relabelSitesApp,
        TNLean.PEPS.EncodedFrame.eval_relabelSitesApp,
        TNLean.PEPS.EncodedFrame.tagSplit,
        TNLean.PEPS.EncodedFrame.tagMerge,
        TNLean.PEPS.EncodedFrame.eval_tagSplit,
        TNLean.PEPS.EncodedFrame.eval_tagMerge,
        TNLean.PEPS.EncodedFrame.exists_tagWord_of_perm}
    \leanok
    \uses{def:peps_frame_registers,lem:peps_frame_registers_coordinates,
        def:peps_monomial_parties}
    The following compositions of exchanges of adjacent tensor factors are
    allowed monomials that use no party and no pair resource, and each sends
    product basis vectors to the indicated product basis vectors: moving the
    raw registers of the sites satisfying a condition in front of the others,
    keeping the order within both parts, and back; reassociating and
    interchanging blocks of registers; splitting the tag registers of a
    concatenated list of holes into those of the two parts, and merging them.
    Naming a register with an equal owner label moves nothing. If a second
    list of holes is a reordering of the first, with pairwise disjoint outer
    footprints, there are a relabelling of the tags that preserves the raw
    part of the encoding and such a composition sending $\ket\tau$ to the
    relabelled $\ket\tau$ (the canonical identification of tag orderings,
    \texttt{05-frames.tex}, lines 84--85).
\end{lemma}

\begin{proof}\leanok
    Each composition is built by induction on the lists, one register or one
    block at a time, and its operator is computed on product vectors. For a
    reordering of the holes, induct on the reordering: an exchange of two
    adjacent holes is an exchange of their tag registers, and the raw parts of
    the encodings of the two holes commute by disjointness of supports.
\end{proof}

\begin{lemma}[Grouped monomials on the registers of a frame]
    \label{lem:peps_grouped_monomial_placement}
    \lean{TNLean.PEPS.EncodedFrame.listT,
        TNLean.PEPS.EncodedFrame.listU,
        TNLean.PEPS.EncodedFrame.listE,
        TNLean.PEPS.EncodedFrame.groupWord,
        TNLean.PEPS.EncodedFrame.ungroupWord,
        TNLean.PEPS.EncodedFrame.eval_groupWord,
        TNLean.PEPS.EncodedFrame.eval_ungroupWord,
        TNLean.PEPS.EncodedFrame.groupWord_props,
        TNLean.PEPS.EncodedFrame.ungroupWord_props,
        TNLean.PEPS.EncodedFrame.merge,
        TNLean.PEPS.EncodedFrame.relabelRest,
        TNLean.PEPS.EncodedFrame.eval_place,
        TNLean.PEPS.EncodedFrame.sheetPlace,
        TNLean.PEPS.EncodedFrame.sum_single_merge,
        TNLean.PEPS.EncodedFrame.layoutIso_place}
    \leanok
    \uses{def:peps_frame_registers,lem:peps_frame_registers_coordinates,
        lem:peps_register_reorderings,lem:peps_front_registers}
    Let $T,E\subseteq\Lambda$ be disjoint and $U=\Lambda\setminus(T\cup E)$,
    let the raw registers of $T$ be held by one party and those of $U$ by
    one party, before a change of raw owners that keeps the owners on $E$,
    and similarly after it. Moving the raw registers of $T$ and of $U$ in
    front by exchanges of tensor factors and grouping each block into one
    register by a unitary at its owner is an allowed monomial using only
    these owners. If a map acts on the two grouped registers, followed by the
    tag registers and the raw registers of $E$, as $M\otimes\id$, then
    grouping at the owners before, applying it, and ungrouping at the owners
    after acts on the registers of the frame, in canonical coordinates, as
    $\id_{\rm tags}\otimes(\id_E\otimes M)$.
\end{lemma}

\begin{proof}\leanok
    On the product basis vector of a configuration $c$, grouping gives
    $\ket{c|_T}\otimes\ket{c|_U}$ in front of the remaining registers; the map
    gives $\sum_{a,b}M_{(a,b),(c|_T,c|_U)}\ket a\otimes\ket b$; ungrouping
    turns each term into the basis vector of the configuration equal to $a$
    on $T$, to $c$ on $E$ and to $b$ on $U$. These are the columns of
    $\id_{\rm tags}\otimes(\id_E\otimes M)$, and both sides are linear.
\end{proof}

\begin{theorem}[Homogeneous birth and death: bounded change]
    \label{thm:peps_frame_birth_bounded}
    \lean{TNLean.PEPS.EncodedFrame.Frame.birth_bounded,
        TNLean.PEPS.EncodedFrame.Frame.death_bounded,
        TNLean.PEPS.PairEffect.PartyChain.precomp,
        TNLean.PEPS.PairEffect.PartyChain.postcomp,
        TNLean.PEPS.PairEffect.PartyChain.eval_precomp,
        TNLean.PEPS.PairEffect.PartyChain.eval_postcomp}
    \leanok
    \uses{thm:peps_frame_birth,def:peps_frame_registers,
        lem:peps_frame_registers_coordinates,def:peps_monomial_parties}
    In an encoded frame $\mathcal F$, let $T$ be a set of raw sites owned by
    $P_\circ$ and disjoint from every outer hole footprint, let $E$ and $U$ be
    as in the birth partition, let $\lVert\Omega\rVert=1$, and suppose
    $I_\Omega(T:E)\le L^{-60}$. Then the birth is a bounded change involving
    only $P_\circ$ and $Q_\circ$, with reference-vector error at most
    $L^{-30}$: there are splitting data and an allowed monomial from the
    registers of $\mathcal F$ to those of the frame in which $T$ is owned by
    $Q_\circ$, using only $P_\circ$ and $Q_\circ$ and at most one normalized
    pair source, whose operator in canonical coordinates is
    $\id_{\rm tags}\otimes B$ for the canonical map $B$, and which maps the
    reference vector of $\mathcal F$ within $L^{-30}$ of the reference vector
    of the new frame. For $Q_\circ\neq P_\circ$ it consists of a private
    contraction at $P_\circ$, one normalized pair source between $Q_\circ$ and
    $P_\circ$, and a private contraction at $P_\circ$, together with
    reorderings and grouping unitaries. The reversed death obeys the same
    bound when these conditions hold in the frame after the death, with one
    normalized pair effect instead and no pair source. This is
    \cite[Lemma~6.5]{OpenAI2026PolynomialPEPS}, \texttt{05-frames.tex},
    lines 396--410.
\end{theorem}

\begin{proof}\leanok
    \uses{thm:peps_frame_birth,lem:peps_frame_birth_partition_owners,
        thm:peps_frame_birth_bounded_grouped,lem:peps_grouped_monomial_placement}
    Take the splitting data of Theorem~\ref{thm:peps_frame_birth}. The
    registers of $T$ and $U$ are held by $P_\circ$ before the birth, and those
    of $T$ by $Q_\circ$ and those of $U$ by $P_\circ$ after it; the owners on
    $E$ are unchanged. Compose the grouping at $P_\circ$, the birth or death
    monomial on grouped registers, and the ungrouping at the owners after the
    change. By Lemma~\ref{lem:peps_grouped_monomial_placement} its operator is
    $\id_{\rm tags}\otimes(\id_E\otimes K^\dagger K)=\id_{\rm
            tags}\otimes B$, and the reference error is that of
    Theorem~\ref{thm:peps_frame_birth}, read through the unitary
    identification of the registers.
\end{proof}

\subsection{Exchanging assignments between two sheets}

\begin{definition}[Swapping two sheets]
    \label{def:peps_frame_sheet_swap}
    \lean{TNLean.PEPS.EncodedFrame.sheetSwap,
        TNLean.PEPS.EncodedFrame.sheetSwapOp}
    \leanok
    For $S\subseteq\Lambda$, $F_S$ is the unitary swapping the two raw sheets
    at the positions in $S$ (\texttt{05-frames.tex}, lines 476--477).
\end{definition}

\begin{lemma}[Composing and commuting sheet swaps]
    \label{lem:peps_frame_sheet_swap_compose}
    \lean{TNLean.PEPS.EncodedFrame.sheetSwapOp_mul_sheetSwapOp,
        TNLean.PEPS.EncodedFrame.commute_sheetSwapOp_kronecker}
    \leanok
    \uses{def:peps_frame_sheet_swap}
    If $Y$ is the disjoint union of $T$ and $A$, then $F_AF_Y=F_T$. The swap
    $F_S$ commutes with $R_1\otimes R_2$ whenever $R_1$ and $R_2$ act on a
    set disjoint from $S$ (\texttt{05-frames.tex}, lines 512--520).
\end{lemma}

\begin{proof}\leanok
    The swaps are permutation matrices of the two-sheet configurations. On a
    configuration, swapping at $Y$ and then at $A\subseteq Y$ swaps exactly
    at the sites of $Y$ outside $A$, that is, at $T$. For the second claim,
    $R_1$ and $R_2$ are the identity on $S$ tensored with operators $R_1'$
    and $R_2'$ on the complement, so in the coordinates splitting off the two
    copies of $S$ both $F_S$ and $R_1\otimes R_2$ act on different factors.
\end{proof}

\begin{definition}[Exchange data]
    \label{def:peps_frame_exchange_data}
    \lean{TNLean.PEPS.EncodedFrame.TwoSheetExchange,
        TNLean.PEPS.EncodedFrame.TwoSheetExchange.frame₁,
        TNLean.PEPS.EncodedFrame.TwoSheetExchange.frame₂,
        TNLean.PEPS.EncodedFrame.TwoSheetExchange.newFrame₁,
        TNLean.PEPS.EncodedFrame.TwoSheetExchange.newFrame₂,
        TNLean.PEPS.EncodedFrame.pairwiseDisjointOuter_append,
        TNLean.PEPS.EncodedFrame.TwoSheetExchange.renameEquiv,
        TNLean.PEPS.EncodedFrame.TwoSheetExchange.rename,
        TNLean.PEPS.EncodedFrame.TwoSheetExchange.newFrame₁_owner_of_mem,
        TNLean.PEPS.EncodedFrame.TwoSheetExchange.newFrame₂_owner_of_mem,
        TNLean.PEPS.EncodedFrame.TwoSheetExchange.ofFrames,
        TNLean.PEPS.EncodedFrame.TwoSheetExchange.mem_footprint_perm,
        TNLean.PEPS.EncodedFrame.TwoSheetExchange.mem_ofFrames_exchangeEnv}
    \leanok
    \uses{def:peps_frame,def:peps_frame_sheet_swap}
    Consider two encoded frames $\mathcal F_1,\mathcal F_2$ and the physical
    sample of a plane region $Y$, such that every hole's outer square lies on
    one side of $\partial Y$; that side classifies the hole as outside or
    inside $Y$. The holes of each frame are listed as those outside $Y$
    followed by those inside $Y$, each in its original order, which differs
    from the given listing by the canonical identification of tag orderings
    (\texttt{05-frames.tex}, lines 84--85). The exchanged
    frames
    $\mathcal F_1',\mathcal F_2'$ carry, inside $Y$, the raw ownership
    assignment and the holes, with their tag owners, of the other frame
    (\texttt{05-frames.tex}, lines 461--464). The renaming $\mathcal R$
    exchanges the sheet names of the raw registers inside $Y$ and of the tags
    of the holes inside $Y$; every register keeps its party.
\end{definition}

\begin{lemma}[Exact intertwining of the renaming]
    \label{lem:peps_frame_exchange_intertwining}
    \lean{TNLean.PEPS.EncodedFrame.TwoSheetExchange.rename_mul_encoder,
        TNLean.PEPS.EncodedFrame.TwoSheetExchange.exists_out_form,
        TNLean.PEPS.EncodedFrame.TwoSheetExchange.exists_in_form,
        TNLean.PEPS.EncodedFrame.sheetSplit_sheetSwap_fst,
        TNLean.PEPS.EncodedFrame.sheetSplit_sheetSwap_snd,
        TNLean.PEPS.EncodedFrame.mul_submatrix_sheetSplit}
    \leanok
    \uses{def:peps_frame_exchange_data,lem:peps_frame_split_forms,
        lem:peps_frame_commute}
    With $K_{\rm in}=K_{\mathcal F_1}\otimes K_{\mathcal F_2}$ and
    $K_{\rm out}=K_{\mathcal F_1'}\otimes K_{\mathcal F_2'}$,
    \begin{align}
        \mathcal RK_{\rm in}=K_{\rm out}F_{Y\cap\Lambda}
    \end{align}
    (\texttt{05-frames.tex}, equation (exchange-intertwining), lines
    490--498).
\end{lemma}

\begin{proof}\leanok
    On every pair of tags the raw part of each encoding is the product of an
    operator acting outside $Y$ and an operator acting on $Y$, which is a
    tensor product in the coordinates $\mathcal H_Y\otimes\mathcal
        H_{\Lambda\setminus Y}$. Swapping the sheets at $Y$ exchanges the
    $Y$-parts, and the four factors match.
\end{proof}

\begin{definition}[Buffer correction]
    \label{def:peps_frame_buffer_correction}
    \lean{TNLean.PEPS.EncodedFrame.tSwap,
        TNLean.PEPS.EncodedFrame.tSwapW,
        TNLean.PEPS.EncodedFrame.bufferSwap,
        TNLean.PEPS.EncodedFrame.tbSwap,
        TNLean.PEPS.EncodedFrame.bufferCorrection,
        TNLean.PEPS.EncodedFrame.threeSplit₂,
        TNLean.PEPS.EncodedFrame.sheetBufferCorrection}
    \leanok
    \uses{lem:peps_frame_splitting}
    With the data of the splitting lemma, the buffer correction on the two
    copies of $U$ is
    \begin{align}
        D_U=(V^{\otimes2})^\dagger F_{\mathcal B_T}V^{\otimes2},
    \end{align}
    where $F_{\mathcal B_T}$ swaps the two copies of $\mathcal B_T$
    (\texttt{05-frames.tex}, equation (exchange-buffer-correction), lines
    503--510). It acts as the identity on the two copies of $T\cup E$.
\end{definition}

\begin{lemma}[Net map of the corrections]
    \label{lem:peps_frame_exchange_net_map}
    \lean{TNLean.PEPS.EncodedFrame.tSwapW_toMatrix,
        TNLean.PEPS.EncodedFrame.tbSwap_toMatrix,
        Matrix.kronecker_kronecker_eq_submatrix,
        Matrix.toMatrix_toPEquiv_prodCongr,
        Matrix.toMatrix_toPEquiv_submatrix,
        TNLean.PEPS.EncodedFrame.bufferCorrection_mul_tSwap,
        TNLean.PEPS.EncodedFrame.sheetSwapOp_eq_submatrix,
        TNLean.PEPS.EncodedFrame.TwoSheetExchange.sheetBufferCorrection_mul_sheetSwapOp}
    \leanok
    \uses{def:peps_frame_buffer_correction,def:peps_frame_sheet_swap}
    With $\mathcal V=\id_{TE}\otimes V$,
    \begin{align}
        D_UF_T=(\mathcal V^{\otimes2})^\dagger F_{T\mathcal B_T}
        \mathcal V^{\otimes2}
    \end{align}
    (\texttt{05-frames.tex}, equation (exchange-net-map), lines 516--526).
\end{lemma}

\begin{proof}\leanok
    In the coordinates in which the two copies of $\mathcal H_T\otimes\mathcal
        H_E$ come first, $\mathcal V^{\otimes2}=\id\otimes V^{\otimes2}$,
    $F_{T\mathcal B_T}=F_T\otimes F_{\mathcal B_T}$, $F_T=F_T\otimes\id$ and
    $D_U=\id\otimes D_U$, so both sides equal $F_T\otimes D_U$.
\end{proof}

\begin{lemma}[Exchange error on the reference]
    \label{lem:peps_frame_exchange_error}
    \lean{EuclideanSpace.vecKron,
        TNLean.PEPS.EncodedFrame.act_kronecker_vecKron,
        EuclideanSpace.norm_vecKron,
        EuclideanSpace.norm_vecKron_self_sub_le,
        TNLean.PEPS.EncodedFrame.act_tbSwap_vecKron_tensorPurification,
        TNLean.PEPS.EncodedFrame.norm_toLp_tensorPurification,
        TNLean.PEPS.EncodedFrame.norm_bufferCorrection_le_one,
        TNLean.PEPS.EncodedFrame.norm_act_bufferCorrection_mul_tSwap_sub_le}
    \leanok
    \uses{lem:peps_frame_exchange_net_map}
    For a unit vector $\omega$ and $\delta=\lVert\mathcal V\omega-s\otimes
        s'\rVert$,
    \begin{align}
        \lVert D_UF_T\,\omega^{\otimes2}-\omega^{\otimes2}\rVert\le4\delta .
    \end{align}
\end{lemma}

\begin{proof}\leanok
    Both $\mathcal V\omega$ and $\zeta=s\otimes s'$ are unit vectors, so
    $\lVert(\mathcal V\omega)^{\otimes2}-\zeta^{\otimes2}\rVert\le2\delta$.
    The swap $F_{T\mathcal B_T}$ exchanges the two identical factors $s$ of
    $\zeta^{\otimes2}$ and fixes it, so by the net map the difference is
    $(\mathcal V^{\otimes2})^\dagger(F_{T\mathcal
            B_T}-\id)\bigl((\mathcal V\omega)^{\otimes2}-\zeta^{\otimes2}\bigr)$,
    of norm at most $2\cdot2\delta$ (\texttt{05-frames.tex}, lines 534--553).
\end{proof}

\begin{definition}[Exchange partition and implemented map]
    \label{def:peps_frame_exchange_map}
    \lean{TNLean.PEPS.EncodedFrame.TwoSheetExchange.exchangeEnv,
        TNLean.PEPS.EncodedFrame.TwoSheetExchange.tSet,
        TNLean.PEPS.EncodedFrame.TwoSheetExchange.eSet,
        TNLean.PEPS.EncodedFrame.TwoSheetExchange.aSet,
        TNLean.PEPS.EncodedFrame.TwoSheetExchange.mem_exchangeEnv,
        TNLean.PEPS.EncodedFrame.liftTags,
        TNLean.PEPS.EncodedFrame.TwoSheetExchange.exchangeOp}
    \leanok
    \uses{def:peps_frame_exchange_data,def:peps_frame_buffer_correction}
    For a party $P_\circ$ put
    \begin{align}
        Z=\{x:\text{the two old raw owners of }x\text{ are not both }P_\circ\}
        \cup H^+(\mathcal F_1)\cup H^+(\mathcal F_2),
    \end{align}
    $T=Z\cap Y$, $E=Z\setminus Y$, $U=\Lambda\setminus Z$ and $A=Y\cap U$
    (\texttt{05-frames.tex}, equations (exchange-Z) and (exchange-partition),
    lines 466--473). The implemented map is
    $C=D_UF_A\mathcal R$, with $D_U$ and $F_A$ acting on raw registers and as
    the identity on all tags.
\end{definition}

\begin{lemma}[Owners outside the exchange environment]
    \label{lem:peps_frame_exchange_owners}
    \lean{TNLean.PEPS.EncodedFrame.TwoSheetExchange.owner_eq_of_notMem_exchangeEnv,
        TNLean.PEPS.EncodedFrame.TwoSheetExchange.newOwner_eq_of_notMem_exchangeEnv}
    \leanok
    \uses{def:peps_frame_exchange_map}
    Every register of $U$ is raw and held by $P_\circ$ on both sheets,
    before and after the exchange (\texttt{05-frames.tex}, lines 500--501).
\end{lemma}

\begin{proof}\leanok
    A register of $U$ is outside $Z$, so both its old raw owners are
    $P_\circ$ and it lies in no outer hole footprint. The exchanged frames
    take their raw owners at each site from one of the two old frames, so
    both new raw owners are $P_\circ$ as well.
\end{proof}

\begin{theorem}[Two-sheet exchange: canonical map]
    \label{thm:peps_frame_exchange}
    \lean{TNLean.PEPS.EncodedFrame.TwoSheetExchange.exchange,
        TNLean.PEPS.EncodedFrame.TwoSheetExchange.exchangeOp_mul_encoder,
        TNLean.PEPS.EncodedFrame.TwoSheetExchange.commute_correction,
        TNLean.PEPS.EncodedFrame.TwoSheetExchange.norm_exchangeOp_le_one,
        TNLean.PEPS.EncodedFrame.liftTags_mul_kronecker,
        TNLean.PEPS.EncodedFrame.commute_sheetBufferCorrection}
    \leanok
    \uses{lem:peps_frame_exchange_intertwining,lem:peps_frame_exchange_net_map,
        lem:peps_frame_exchange_error,def:peps_frame_exchange_map,
        lem:peps_frame_sheet_swap_compose,lem:peps_frame_exchange_owners}
    In the setting of the exchange partition, let $\lVert\Omega\rVert=1$. If
    $I_\Omega(T:E)\le L^{-60}$, then for splitting data with error at most
    $L^{-30}$ the map $C$ is a contraction,
    $CK_{\rm in}=K_{\rm out}D_UF_T$, and
    \begin{align}
        \lVert C(\Omega_{\mathcal F_1}\otimes\Omega_{\mathcal F_2})
        -\Omega_{\mathcal F_1'}\otimes\Omega_{\mathcal F_2'}\rVert
        \le4L^{-30}.
    \end{align}
    The encoded reference vectors are not assumed normalized. This is the
    content of \cite[Lemma~6.6]{OpenAI2026PolynomialPEPS} for the implemented
    map, with its hypotheses, when the holes of each frame are listed as those
    outside $Y$ followed by those inside $Y$.
\end{theorem}

\begin{proof}\leanok
    \uses{lem:peps_frame_splitting,lem:peps_frame_exchange_intertwining,
        def:peps_frame_sheet_swap,lem:peps_frame_sheet_swap_compose,
        lem:peps_frame_exchange_net_map,lem:peps_frame_exchange_error,
        def:peps_frame_exchange_map,lem:peps_frame_exchange_owners,
        lem:peps_frame_encoder_contraction}
    The corrections act on $U$, away from every hole of both sheets, so they
    commute with $K_{\rm out}$; with $F_AF_Y=F_T$ and the intertwining of the
    renaming, $CK_{\rm in}=K_{\rm out}D_UF_T$. The encoding $K_{\rm out}$ is a
    contraction, and the exchange error on the reference with
    $\delta\le L^{-30}$ gives the bound.
\end{proof}

\begin{theorem}[Two-sheet exchange: holes listed in any order]
    \label{thm:peps_frame_exchange_any_order}
    \lean{TNLean.PEPS.EncodedFrame.TwoSheetExchange.exchange_ofFrames,
        TNLean.PEPS.EncodedFrame.TwoSheetExchange.tagReorder,
        TNLean.PEPS.EncodedFrame.TwoSheetExchange.tagReorder_mul_kronecker,
        TNLean.PEPS.EncodedFrame.TwoSheetExchange.norm_tagReorder_le_one}
    \leanok
    \uses{thm:peps_frame_exchange}
    Let $\mathcal F_1,\mathcal F_2$ be encoded frames whose holes are listed in
    any order, each hole classified as outside or inside $Y$ as in the
    exchange data, and let $\lVert\Omega\rVert=1$ and $I_\Omega(T:E)\le
        L^{-60}$. There are a relabelling $G$ of the tags of both frames with
    $G(K_{\mathcal F_1}\otimes K_{\mathcal F_2})=K_{\rm in}$, the encodings
    of the regrouped listings, and splitting data with error at most
    $L^{-30}$, such that $CG$ is a contraction,
    $CG(K_{\mathcal F_1}\otimes K_{\mathcal F_2})=K_{\rm out}D_UF_T$, and
    \begin{align}
        \lVert CG(\Omega_{\mathcal F_1}\otimes\Omega_{\mathcal F_2})
        -\Omega_{\mathcal F_1'}\otimes\Omega_{\mathcal F_2'}\rVert
        \le4L^{-30}.
    \end{align}
\end{theorem}

\begin{proof}\leanok
    \uses{thm:peps_frame_exchange,thm:peps_frame_small_rewrite,
        lem:peps_frame_exchange_error}
    Encodings of distinct holes commute after the canonical identification of
    tag orderings (\texttt{05-frames.tex}, lines 84--85), which gives $G$.
    Apply the exchange to the regrouped listings and compose with $G$, a
    contraction.
\end{proof}

\begin{lemma}[Exchange corrections as one private contraction]
    \label{lem:peps_frame_exchange_correction_monomial}
    \lean{TNLean.PEPS.EncodedFrame.partSwap,
        TNLean.PEPS.EncodedFrame.correctionMatrix,
        TNLean.PEPS.EncodedFrame.norm_correctionMatrix_le_one,
        TNLean.PEPS.EncodedFrame.sheetSwapOp_eq_submatrix_of_subset,
        TNLean.PEPS.EncodedFrame.sheetBufferCorrection_mul_sheetSwapOp_eq,
        TNLean.PEPS.EncodedFrame.correctionWord,
        TNLean.PEPS.EncodedFrame.TwoSheetExchange.exchange_correction_monomial}
    \leanok
    \uses{def:peps_frame_exchange_map,lem:peps_frame_exchange_owners,
        lem:peps_front_registers,def:peps_monomial_parties}
    In the setting of the exchange partition, the corrections are
    $D_UF_A=\id\otimes W$ in the coordinates in which the two copies of
    $\mathcal H_T\otimes\mathcal H_E$ come first, where
    $W=(V^{\otimes2})^\dagger F_{\mathcal B_T}V^{\otimes2}F_A$ acts on the two
    copies of $\mathcal H_U$ and $\lVert W\rVert\le1$. The registers of $U$ are
    held by $P_\circ$ on both sheets before and after the exchange. On two
    registers of $P_\circ$, each grouping the raw registers of $U$ on one sheet,
    followed by an arbitrary list of further registers, $W\otimes\id$ is the
    operator of one allowed private contraction at $P_\circ$
    (\texttt{05-frames.tex}, lines 503--515).
\end{lemma}

\begin{proof}\leanok
    Since $A\subseteq U$, the swap $F_A$ exchanges only the $A$-parts of the
    two copies of $\mathcal H_U$, so $F_A=\id\otimes F'_A$, and
    $D_U=\id\otimes(V^{\otimes2})^\dagger F_{\mathcal B_T}V^{\otimes2}$. Each
    factor of $W$ is a contraction.
\end{proof}

\begin{lemma}[Registers of two frames]
    \label{lem:peps_two_frame_registers}
    \lean{TNLean.PEPS.EncodedFrame.twoLayoutIso,
        TNLean.PEPS.EncodedFrame.twoLayoutIso_symm_single,
        TNLean.PEPS.EncodedFrame.twoLayoutIso_basis,
        TNLean.PEPS.EncodedFrame.twoLayoutIso_eq_act_of_basis,
        TNLean.PEPS.EncodedFrame.act_toMatrix_symm_single}
    \leanok
    \uses{lem:peps_frame_registers_coordinates}
    The registers of two frames, those of the first followed by those of the
    second, are identified with the product of the canonical coordinates of
    the two frames, product basis vectors going to standard basis vectors. Two
    operators on these registers agree as soon as they agree on product basis
    vectors.
\end{lemma}

\begin{proof}\leanok
    Combine the identifications of the two frames, and expand a vector in the
    standard basis.
\end{proof}

\begin{lemma}[The renaming as reorderings]
    \label{lem:peps_exchange_renaming_word}
    \lean{TNLean.PEPS.EncodedFrame.TwoSheetExchange.sitesY,
        TNLean.PEPS.EncodedFrame.TwoSheetExchange.sitesN,
        TNLean.PEPS.EncodedFrame.TwoSheetExchange.renameWord,
        TNLean.PEPS.EncodedFrame.TwoSheetExchange.eval_renameWord,
        TNLean.PEPS.EncodedFrame.TwoSheetExchange.twoLayoutIso_renameWord,
        TNLean.PEPS.EncodedFrame.TwoSheetExchange.renameWord_props}
    \leanok
    \uses{def:peps_frame_exchange_data,lem:peps_register_reorderings,
        lem:peps_two_frame_registers}
    The renaming $\mathcal R$ is the operator of a composition of exchanges
    of tensor factors from the registers of the two frames to the registers
    of the two exchanged frames, using no party and no pair resource. It
    splits the tags of each frame into those of the holes outside and inside
    $Y$, moves the raw registers of the sites of $Y$ next to the inside tags,
    exchanges these blocks between the sheets, and merges back; every
    register keeps its party (\texttt{05-frames.tex}, lines 485--490).
\end{lemma}

\begin{proof}\leanok
    Follow a product basis vector through the reorderings. Inside $Y$ the
    configurations and the inside tags of the two sheets change places, and
    the owners of the moved raw registers are the new owners at those sites,
    so the result is the basis vector of the renamed configuration.
\end{proof}

\begin{lemma}[The corrections on the registers of two frames]
    \label{lem:peps_exchange_correction_word}
    \lean{TNLean.PEPS.EncodedFrame.sitesU,
        TNLean.PEPS.EncodedFrame.sitesR,
        TNLean.PEPS.EncodedFrame.pairUIso,
        TNLean.PEPS.EncodedFrame.correctionLayoutWord,
        TNLean.PEPS.EncodedFrame.eval_correctionLayoutWord,
        TNLean.PEPS.EncodedFrame.twoSheetPlace,
        TNLean.PEPS.EncodedFrame.sum_single_merge₂,
        TNLean.PEPS.EncodedFrame.twoLayoutIso_correctionLayoutWord,
        TNLean.PEPS.EncodedFrame.correctionLayoutWord_props}
    \leanok
    \uses{lem:peps_register_reorderings,lem:peps_two_frame_registers,
        lem:peps_frame_exchange_correction_monomial}
    Let $T,E\subseteq\Lambda$ be disjoint and let the raw registers of
    $U=\Lambda\setminus(T\cup E)$ be held by $P_\circ$ on both frames. Moving
    the raw registers of $U$ of both sheets to the front, applying a
    contraction $W$ of the two copies of $\mathcal H_U$ as one private
    contraction at $P_\circ$, and moving them back is an allowed monomial
    using only $P_\circ$ and no pair resource, whose operator in canonical
    coordinates is the identity on the tags and on the two copies of
    $T\cup E$, tensored with $W$.
\end{lemma}

\begin{proof}\leanok
    As for the grouped monomials on one frame: on a product basis vector the
    moved registers carry the two configurations on $U$, the contraction
    gives the column of $W$, and moving back gives the basis vectors with
    these configurations replaced.
\end{proof}

\begin{theorem}[Two-sheet exchange: bounded change]
    \label{thm:peps_frame_exchange_bounded}
    \lean{TNLean.PEPS.EncodedFrame.TwoSheetExchange.exchange_bounded,
        TNLean.PEPS.EncodedFrame.TwoSheetExchange.exchange_ofFrames_bounded,
        TNLean.PEPS.EncodedFrame.exists_tagReorderWord}
    \leanok
    \uses{thm:peps_frame_exchange,thm:peps_frame_exchange_any_order,
        def:peps_frame_registers,def:peps_monomial_parties}
    Consider two encoded frames and a region $Y$ such that every hole's
    outer square lies on one side of $\partial Y$, a party $P_\circ$ with
    $Z$, $T$, $E$, $U$ as in the exchange partition, and
    $\lVert\Omega\rVert=1$ with $I_\Omega(T:E)\le L^{-60}$. The exchange is
    implemented by private contractions and register renaming, with
    reference-vector error at most $4L^{-30}$: there are splitting data and
    an allowed monomial from the registers of the two frames to those of the
    two exchanged frames, consisting of exchanges of tensor factors and one
    private contraction at $P_\circ$, using only $P_\circ$ and no pair
    resource, whose operator in canonical coordinates is the implemented map
    $C=D_UF_A\mathcal R$ (composed with the relabelling $G$ of the tags when
    the holes are listed in any order), and which maps the product of the two
    reference vectors within $4L^{-30}$ of the product of the new ones. In
    particular the exchange is a bounded change in the sense of
    \cite[Theorem~5.2]{OpenAI2026PolynomialPEPS}. This is
    \cite[Lemma~6.6]{OpenAI2026PolynomialPEPS}, \texttt{05-frames.tex},
    lines 460--479; the source allows also the owners of the renamed registers
    and tags, which exchanges of tensor factors do not involve.
\end{theorem}

\begin{proof}\leanok
    \uses{thm:peps_frame_exchange,thm:peps_frame_exchange_any_order,
        lem:peps_exchange_renaming_word,lem:peps_exchange_correction_word,
        lem:peps_frame_exchange_correction_monomial,
        lem:peps_frame_exchange_owners,lem:peps_register_reorderings}
    Compose the renaming as reorderings with the corrections on the
    registers of the exchanged frames, whose raw registers of $U$ are held by
    $P_\circ$ on both sheets. The operators compose to
    $D_UF_A\mathcal R$, since $D_UF_A$ is the identity on the tags and on the
    two copies of $T\cup E$ tensored with one contraction of the two copies
    of $\mathcal H_U$. For holes in any order, precede by the canonical
    identification of tag orderings as reorderings, whose operator is $G$.
    The reference error is that of Theorems~\ref{thm:peps_frame_exchange}
    and~\ref{thm:peps_frame_exchange_any_order}, read through the unitary
    identification of the registers.
\end{proof}
